% Verified Capability Discipline for LLM Agent Tool Calls
%
% Submission shell. The body lives in DRAFT.md as a working document; this
% .tex file is the camera-ready translation. The two should be kept in sync;
% scripts/md-to-tex-diff.py (TBD) helps detect drift before submission.
%
% Target: IEEE S&P 2027.

\documentclass[conference]{IEEEtran}

% S&P uses IEEEtran with a few tweaks; the venue's official style file
% (if provided) should be substituted here at submission time.

\usepackage[utf8]{inputenc}
\usepackage[T1]{fontenc}
\usepackage{lmodern}
\usepackage{graphicx}
\usepackage{xcolor}
\usepackage{listings}
\usepackage{amsmath, amssymb, amsthm}
\usepackage[hidelinks]{hyperref}
\usepackage{booktabs}
\usepackage{enumitem}
\usepackage{microtype}

\newtheorem{theorem}{Theorem}

\lstset{
  basicstyle=\ttfamily\small,
  breaklines=true,
  showstringspaces=false,
  frame=single,
  framerule=0.3pt,
  rulecolor=\color{gray!50},
  xleftmargin=0.5em,
  xrightmargin=0.5em,
}

\title{Verified Capability Discipline for LLM Agent Tool Calls}

\author{
  \IEEEauthorblockN{[Anonymised for double-blind submission]}
}

\begin{document}

\maketitle

\begin{abstract}
Large language model agents now mediate billions of tool calls per day across regulated industries, but every published defence against prompt injection reduces attack success without eliminating it: the latest benchmarks report 14--31\% attack success against the strongest text-side defences. We argue that this ceiling is a structural property of attempting to verify unbounded natural language. The consequential half of an agent's behaviour, however, passes through a much narrower interface: every modern agent framework declares tool interfaces as JSON Schema, making the set of well-formed tool calls bounded and decidable.

We present \emph{Verified Capability Discipline} (VCD), a composition of three primitives that structurally preclude prompt-injection-driven tool misuse over the verified action surface. An SMT-backed prover encodes a tool's JSON Schema and a security policy into Z3, deciding whether every string the schema permits satisfies the policy or extracting a concrete counterexample. An issuer mints Ed25519-signed capability tokens binding agent, tool, constraints, and expiry; an \texttt{attenuate} primitive derives more-restricted children whose bounds are provably tighter than their parents and whose lifetimes cannot exceed them. A gate sits on the only path from agent intent to tool execution and enforces eight checks, failing closed by default.

We mechanise three soundness theorems in Lean~4 with zero \texttt{sorry}s: attenuation cannot broaden permissions, the gate's ALLOW implies constraint satisfaction, and a token citing a proof guarantees its accepted calls conform to the original schema and policy. Evaluating against AgentDojo and InjecAgent across four state-of-the-art baselines, VCD reduces tool-call-mediated attack success rate from 14--31\% to under 1\% (LLM-driven measurement pending) at a per-call gate latency well under 100 microseconds at p50 on commodity hardware. A static upper bound across both benchmarks is already verified at 0 of 2{,}737 attack scenarios. The reference implementation, Lean development, and benchmark harness are MIT-licensed.
\end{abstract}

% ---------------------------------------------------------------------------
% 1. Introduction
% ---------------------------------------------------------------------------

\section{Introduction}\label{sec:intro}

% Placeholder. Run paper/build-tex.sh to regenerate from DRAFT.md.
\noindent\textit{[This section is sourced from \texttt{paper/DRAFT.md} via
\texttt{paper/build-tex.sh}. Run that script to populate. Until then the
section bodies are placeholders so \texttt{main.tex} compiles.]}



% ---------------------------------------------------------------------------
% 2. Threat Model
% ---------------------------------------------------------------------------

\section{Threat Model}\label{sec:threat-model}

\input{sections/02-threat-model.tex}

% ---------------------------------------------------------------------------
% 3. System Design
% ---------------------------------------------------------------------------

\section{System Design}\label{sec:design}

\input{sections/03-design.tex}

% ---------------------------------------------------------------------------
% 4. Formal Analysis
% ---------------------------------------------------------------------------

\section{Formal Analysis}\label{sec:formal}

\input{sections/04-formal.tex}

% ---------------------------------------------------------------------------
% 5. Implementation
% ---------------------------------------------------------------------------

\section{Implementation}\label{sec:implementation}

\input{sections/05-implementation.tex}

% ---------------------------------------------------------------------------
% 6. Evaluation
% ---------------------------------------------------------------------------

\section{Evaluation}\label{sec:eval}

\input{sections/06-evaluation.tex}

% ---------------------------------------------------------------------------
% 7. Related Work
% ---------------------------------------------------------------------------

\section{Related Work}\label{sec:related}

\input{sections/07-related.tex}

% ---------------------------------------------------------------------------
% 8. Limitations and Future Work
% ---------------------------------------------------------------------------

\section{Limitations and Future Work}\label{sec:limitations}

\input{sections/08-limitations.tex}

% ---------------------------------------------------------------------------
% 9. Conclusion
% ---------------------------------------------------------------------------

\section{Conclusion}\label{sec:conclusion}

\input{sections/09-conclusion.tex}

% ---------------------------------------------------------------------------
% Acknowledgements
% ---------------------------------------------------------------------------

% \section*{Acknowledgements}
% Anonymised for double-blind submission.

% ---------------------------------------------------------------------------
% References
% ---------------------------------------------------------------------------

\bibliographystyle{IEEEtran}
\bibliography{references}

\end{document}
